\section{Addressed compact curls}
\label{sec:addressed}

The first construction assigns a localized curl to each point of a fine
spatial grid. Its displacement depends on the configuration encoded by
that address. The curl carries its center along the prescribed path,
while disjoint supports keep neighboring instructions separate. Retaining
old configuration codes at coarser scales lets these small motions
implement transitions that need not be reversible.

\begin{theorem}[Compact support and rapid decay]
\label{thm:rapid}
Fix a positive computable viscosity $\nu$. From a deterministic Turing machine $M$ and a
finite input $w$ one can effectively construct a smooth force on
$[0,\infty)\times\R^3$, supported for all times in the fixed compact set
\[
 K=[-2,1]\times[-1,2]\times[-1,1],
\]
such that, for every pair of nonnegative integers $J,l$ and spatial
multi-index $\alpha$,
\begin{equation}
 \sup_{t\geq0,\,x\in\R^3}
 (1+t)^J\bigl|\partial_t^l\partial_x^\alpha f(t,x)\bigr|<\infty.
 \label{eq:rapid-force-bound}
\end{equation}
The Navier--Stokes equation with this force and zero initial velocity has
an explicit global smooth solution, unique in the finite-interval class \eqref{eq:energy-class}.  Its material trajectory from
the origin satisfies
\begin{equation}
 M\text{ halts on }w
 \quad\Longleftrightarrow\quad
 X_1(t,0)<-1\text{ for some }t\geq0.
 \label{eq:rapid-event}
\end{equation}
The velocity has the same fixed support and rapid mixed-derivative bounds.
The pressure is zero. The force is a general body force. Because $K$ is fixed,
the same bounds hold with $(1+t+|x|)^J$ in place of $(1+t)^J$.
\end{theorem}

\subsection{One transition on a complete finite configuration}
Use a tape indexed by $\mathbb Z$, a head initially at $0$, and input symbols
in cells $0,\ldots,B-1$, where $B=|w|$, with all other cells blank.
A transition writes one symbol and
moves the head by $d\in\{-1,0,1\}$.  Normalize missing transitions to enter
a fresh terminal state without moving or changing the scanned symbol,
and give terminal states no outgoing transitions.  This preserves halting.
Write $Q$, $Q_h\subset Q$, and $\Gamma$ for the resulting state set, terminal
set, and alphabet.  Number the states and symbols consecutively from
zero, giving the blank symbol number zero.  Choose integers $r_A,r_Q\geq1$
so that
\[
 A=2^{r_A}\geq|\Gamma|,\qquad S=2^{r_Q}\geq|Q|.
\]
After $n$ transitions, the head $i$ satisfies $|i|\leq n$, and every
nonblank cell lies in $[-n,n+B]$.  Define
\begin{align}
 P_n&=2^{n+2},&
 M_n&=S P_n A^{2n+B+1}=2^{L_n},\\
 &&
 L_n&=r_Q+n+2+r_A(2n+B+1). \label{eq:rapid-capacity}
\end{align}
For the configuration $(q,\tau,i)$ at that time, put
\begin{equation}
 T=\sum_{v=-n}^{n+B}\operatorname{id}(\tau(v))A^{v+n},\qquad
 j=i+n,\qquad
 k_n=\operatorname{id}(q)+S(j+P_nT).
 \label{eq:rapid-code}
\end{equation}
Here $0\leq j\leq2n<P_n$, so $0\leq k_n<M_n$.  The code and $n$ recover
the state, the head, and the entire tape, with blank extensions.

The following functions are defined on \emph{every} integer
$0\leq k<M_n$, not only on configurations reached from the input.  Extract
\[
 r=k\bmod S,\qquad
 j=\lfloor k/S\rfloor\bmod P_n,\qquad
 T=\lfloor k/(SP_n)\rfloor.
\]
Let $h_n(k)$ be one if $r$ identifies a terminal state and zero otherwise.
If $r$ identifies a nonterminal state, $j\leq2n$, and
$z=\lfloor T/A^j\rfloor\bmod A$ identifies an alphabet symbol, look up the
single transition $(q,z)\mapsto(q',z',d)$ and set
\begin{equation}
 F_n(k)=\operatorname{id}(q')+
 S\bigl(j+d+1+P_{n+1}A(T+(z'-z)A^j)\bigr).
 \label{eq:rapid-transition}
\end{equation}
In all other cases set $F_n(k)=0$.  In the lookup case, replacing the
$j$th digit gives an integer in $[0,A^{2n+B+1})$.  Multiplication by $A$
adds the new blank cell on the left, and the two-place enlargement of the
tape window leaves a blank cell on the right.  Also
$0\leq j+d+1\leq2n+2<P_{n+1}$.  Consequently
\begin{equation}
 0\leq F_n(k)<M_{n+1},
 \qquad k_{n+1}=F_n(k_n)
 \quad\hbox{at each actual nonterminal step}.
 \label{eq:rapid-transition-range}
\end{equation}
The range assertion does not require unscanned digits of a padded code
to represent genuine alphabet symbols.

\subsection{Reading the newest record while retaining the earlier ones}
Set
\begin{equation}
 \begin{gathered}
 C=\max\{1,L_0,1+2r_A\},\qquad s_n=C2^{(n+2)^2},\\
 \delta_n=2^{-s_n},\qquad D_n=2^{s_n},\qquad
 b_n=\delta_{n+1}M_{n+1}.
 \end{gathered}
 \label{eq:rapid-scales}
\end{equation}
These quantities depend on the machine and input length, not on its
execution. The elementary bounds are
\[
 L_n\leq C(n+1),\qquad
 s_n\geq16C(n+1),\qquad
 \frac{s_{n+1}}{s_n}=2^{2n+5}\geq32(n+1).
\]
These estimates give
\begin{align}
 &\sum_{n\geq0}\delta_n M_n
       \leq\sum_{n\geq0}2^{-15(n+1)}<1,
       \label{eq:rapid-total-memory}\\
 &b_n\leq\delta_n^{n+10}<\delta_n/8,
       \label{eq:rapid-small-motion}\\
 &0\leq j<n\quad\Longrightarrow\quad
       \delta_j/\delta_n\text{ is an integer multiple of }M_n.
       \label{eq:rapid-divisibility}
\end{align}
For the second assertion, use
$s_{n+1}-L_{n+1}\geq(15/16)s_{n+1}
\geq30(n+1)s_n\geq(n+10)s_n$.
For the third, $s_n-L_n\geq(15/16)s_n\geq30s_{n-1}\geq s_j$
when $j<n$, using $s_n/s_{n-1}\geq32$.
The strict inequality in \eqref{eq:rapid-small-motion} follows from
$\delta_n\leq2^{-16}$.

Up to the first halt, store the actual configurations at
\begin{equation}
 y_n=\sum_{j=0}^n\delta_j k_j\in[0,1).
 \label{eq:rapid-history}
\end{equation}
This number is an integer multiple $m\delta_n$, with $0\leq m<D_n$.
Equation~\eqref{eq:rapid-divisibility} gives the essential read operation:
\begin{equation}
 m\bmod M_n=k_n.
 \label{eq:rapid-read}
\end{equation}
Thus the current code can be read without discarding the old codes.
This uses the information-retention principle of reversible computation
\cite{Landauer,Bennett}, with spatial digits in place of a history tape.
Retaining this history permits a diffeomorphic flow to realize a machine
whose transition map is not injective: in the realization below, distinct
grid addresses start at least $\delta_n$ apart in the coding coordinate
and move by less than $\delta_n/8$, even when their next codes coincide.

\subsection{Moving each address by a localized curl}
With the switch $H$ and progress function $\theta$ of Section~\ref{sec:analytic}, set
\[
 \zeta(v)=\prod_{\ell=1}^3
   H\bigl(16(v_\ell+1/8)\bigr)
   H\bigl(16(1/8-v_\ell)\bigr).
\]
The nondecreasing function $\theta$ is zero for $s\leq1/4$ and one for
$s\geq3/4$.  The bump $\zeta$ is supported in $[-1/8,1/8]^3$ and equals
one on $[-1/16,1/16]^3$.  All required derivative bounds are effective.

For a prescribed center path $c(\sigma)$ and a spatial scale $\delta>0$,
the curl field
\begin{equation}
 \mathcal V(\sigma,x)=
 \nabla_x\times\left[
  \zeta\left(\frac{x-c(\sigma)}{\delta}\right)
  \frac{\dot c(\sigma)\times(x-c(\sigma))}{2}
 \right]
 \label{eq:rapid-curl}
\end{equation}
is divergence-free and equals $\dot c(\sigma)$ on a neighborhood of the
moving center.  Indeed, there the cutoff is one and
$\nabla_x\times(a\times(x-c))=2a$ for constant $a$.

On $0\leq t\leq1$, use \eqref{eq:rapid-curl} with
$\delta=1$ and $c(t)=\theta(t)(-1,y_0,0)$, where only the initial code
$k_0$ is computed from the input.  This defines a velocity $W$ carrying
the origin exactly to $(-1,y_0,0)$.

On the $n$th subsequent interval $[1+n,2+n]$, let $\sigma=t-1-n$.
For each integer $0\leq m\leq D_n$ set
\begin{align}
 k&=m\bmod M_n,&
 a_{n,m}&=(-1,m\delta_n,0),\\
 d_{n,m}&=\bigl(-b_n h_n(k),\,\delta_{n+1}F_n(k),\,0\bigr),&
 c_{n,m}(\sigma)&=a_{n,m}+\theta(\sigma)d_{n,m}.
 \label{eq:rapid-address-motion}
\end{align}
Let $V_{n,m}$ be \eqref{eq:rapid-curl} with center $c_{n,m}$ and scale
$\delta_n$, and define
\begin{equation}
 W(t,x)=\sum_{m=0}^{D_n}V_{n,m}(\sigma,x)
 \qquad(1+n\leq t\leq2+n).
 \label{eq:rapid-field}
\end{equation}
By \eqref{eq:rapid-transition-range} and \eqref{eq:rapid-small-motion},
$|d_{n,m}|_\infty\leq b_n<\delta_n/8$.  The support of $V_{n,m}$ is
therefore contained in
\[
 |x_2-m\delta_n|<\delta_n/4.
\]
These supports are pairwise disjoint.  In particular the sum can also be
evaluated by the nearest grid index; it vanishes on a neighborhood of
each boundary between two such choices.  Each summand is a smooth curl,
and all time pieces vanish on fixed-width neighborhoods of their joins.
Every finite time interval meets only finitely many pieces, each with
finitely many summands.  Hence $W$ is globally smooth and divergence-free,
with $W(0)=0$.  The center and displacement bounds, including the loading
segment, place its support in the interior of $K$ for all time.

\subsection{Why every fixed derivative decreases rapidly}
For one summand introduce
\[
 \bar x=\frac{x-a_{n,m}}{\delta_n},\qquad
 \eta=\frac{d_{n,m}}{\delta_n},\qquad
 \xi=\frac{d_{n,m}}{b_n}.
\]
Then $V_{n,m}$ is $b_n$ times the fixed smooth expression
\begin{equation}
 \nabla_{\bar x}\times\left[
  \zeta(\bar x-\theta(\sigma)\eta)
  \frac{\theta'(\sigma)\xi\times
          (\bar x-\theta(\sigma)\eta)}{2}
 \right].
 \label{eq:rapid-scaled-field}
\end{equation}
The parameters $\eta,\xi$ lie in $[-1,1]^3$.  On this compact parameter
set the expression has bounded $\sigma,\bar x$ derivatives of every order,
uniformly bounded spatial support, and effective bounds.  Disjointness
of the summands therefore gives, independently of $n$ and $m$,
\begin{equation}
 \|\partial_t^l\partial_x^\alpha W(t,\cdot)\|_\infty
 \leq C_{l,\alpha}b_n\delta_n^{-|\alpha|}
 \qquad(1+n\leq t\leq2+n).
 \label{eq:rapid-velocity-estimate}
\end{equation}
In particular there is no factor counting the grid cells.  Time
differentiation costs no additional negative power of $\delta_n$:
the center displacement in scaled coordinates is the bounded parameter
$\eta$ in \eqref{eq:rapid-scaled-field}.

Prescribe the force
\begin{equation}
 f=\partial_t W+(W\cdot\nabla)W-\nu\Delta W.
 \label{eq:rapid-residual}
\end{equation}
The product rule and \eqref{eq:rapid-velocity-estimate} give on the $n$th
step interval
\begin{equation}
 \|\partial_t^l\partial_x^\alpha f(t,\cdot)\|_\infty
 \leq C'_{l,\alpha,\nu}
 \left(b_n\delta_n^{-|\alpha|-2}
             +b_n^2\delta_n^{-|\alpha|-1}\right).
 \label{eq:rapid-residual-estimate}
\end{equation}
For fixed derivative orders and all sufficiently large $n$,
\eqref{eq:rapid-small-motion} bounds the right-hand side by a constant
times $\delta_n$. More explicitly, its two terms are bounded by
$\delta_n^{n+8-|\alpha|}$ and $\delta_n^{2n+19-|\alpha|}$;
$n\geq|\alpha|$ is a sufficient effective cutoff for both exponents
to be at least one. Since $\delta_n\leq2^{-n}$ and $1+t\leq n+3$ on
this interval, multiplication by any fixed power of $1+t$ remains
bounded.  The finitely many preceding intervals and the loading interval
have finite bounds individually.  This proves
\eqref{eq:rapid-force-bound}, and the same argument applies to $W$.
All constants needed for a prescribed finite set of orders can be
computed: the tail uses the displayed inequalities, and the remaining
finite set uses derivatives of the fixed cutoffs.

Lemma~\ref{lem:realization} now applies to $u=W$, $p=0$, with the force
\eqref{eq:rapid-residual}.  Compact support and the derivative bounds
supply every required finite-interval Sobolev and boundedness condition.
It gives the asserted global solution, uniqueness, and well-defined
material flow.

\subsection{The particle reads its own current configuration}
At time $1$, the observed particle is $(-1,y_0,0)$, and during loading
its first coordinate is $-\theta(t)\geq-1$.  Suppose the machine has
reached its $n$th configuration and the particle is $(-1,y_n,0)$ at
time $1+n$.  By \eqref{eq:rapid-read}, this point is the center
$a_{n,m}$ for $m=y_n/\delta_n$, whose residue is $k_n$.  The curve
$c_{n,m}(\sigma)$ solves the particle equation by
\eqref{eq:rapid-curl} and the disjoint supports, so uniqueness identifies
it with the actual trajectory throughout the interval.

If the configuration is nonterminal, its first coordinate remains
exactly $-1$ throughout that interval.  Its second-coordinate endpoint
is $y_n+\delta_{n+1}F_n(k_n)=y_{n+1}$.  This proves the configuration
correspondence by induction, until a first halt or indefinitely.
If the configuration is terminal, the next interval instead ends with
first coordinate $-1-b_n<-1$.  This also covers an initially terminal
machine.  Thus a halt produces the event in finite time, whereas a
nonhalting run never produces it, even between the coding times.

\subsection{Finite evaluation without executing the input}
Equations~\eqref{eq:rapid-capacity}--\eqref{eq:rapid-field} are a finite
program for the prescribed field.  At any finite time only a bounded
number of possible time pieces need be considered.  Their parameters
are finite integers and dyadic rationals, although no efficiency is
claimed.  Each piece is a finite sum of effective smooth functions.
Equivalently, a second-coordinate approximation with error less than
$\delta_n/32$ locates the only possibly active address unless its error
interval meets a midpoint between adjacent grid points.  In that case
the true coordinate is within $\delta_n/16$ of the midpoint, and hence
outside every supporting strip of half-width $\delta_n/4$; the value
and all derivatives are zero. An inferred grid index outside
$[0,D_n]$ contributes zero. A time approximation with certified error
less than $1/32$ similarly resolves the time piece, or, if its enclosure
meets an integer join, places the true time within $1/16$ of that join.
This lies inside the quarter-unit zero collar. Thus the field, its derivatives, and
\eqref{eq:rapid-residual} can be evaluated to arbitrary precision,
including at the joins, without deciding equality of a real number to
a branch boundary.

Crucially, $F_n$ performs one transition lookup on the integer supplied
by a \emph{spatial address}.  Its definition applies simultaneously to
all padded configurations and never computes $k_1,k_2,\ldots$ from
$k_0$.  Only $k_0$ is used in loading; the particle's successive positions
perform the iteration.  After loading, the entire prescription depends
on the input only through its length.  Physical steps occupy successive
unit intervals, so there is no finite accumulation of computation times.
This completes the proof of Theorem~\ref{thm:rapid}.

\subsection{A solenoidal force on the unit torus}

The compact curl field also fits into a fixed interior torus chart.
The scale change is made at the velocity level; the force is then recomputed
for the specified viscosity.

\begin{corollary}[Solenoidal rapid forcing on the torus]
\label{cor:rapid-torus}
Fix a positive computable viscosity.  From any deterministic Turing
machine $M$ and finite input $w$, one can effectively construct on the
unit flat torus $\T^3$ a smooth, solenoidal, spatially mean-zero
force satisfying all the bounds \eqref{eq:rapid-force-bound} with $\T^3$
in place of $\R^3$, with zero
initial velocity, such that the unique smooth solution satisfies
\[
 M\text{ halts on }w
 \quad\Longleftrightarrow\quad
 \frac14<X_1(t,a_*)<\frac38\text{ for some }t\geq0,
 \qquad a_*=(1/2,1/4,1/2).
\]
Here the inequality denotes the indicated open coordinate slab of the
torus. Pressure is normalized to mean zero, and uniqueness holds in the
classical periodic class of Section~\ref{sec:analytic}.  The velocity and all its mixed derivatives also decrease faster
than every inverse power of time.
\end{corollary}

\begin{proof}
Let $W$ be the divergence-free velocity constructed in
Theorem~\ref{thm:rapid}; its definition does not involve viscosity.
Put $\lambda=1/8$ and $c=(1/2,1/4,1/2)$, and define in the unit chart
\[
 v(t,x)=\lambda W\left(t,\frac{x-c}{\lambda}\right),
\]
extending by zero and then periodically.  The fixed support $c+\lambda K$
lies strictly inside the unit cube.  Thus the extension is smooth,
divergence-free, zero initially, and rapidly decreasing with every mixed
derivative. The velocity $W$ is a curl at each time, with the compactly
supported potentials displayed in \eqref{eq:rapid-curl}.
If $W=\nabla_y\times A$, then $v$ is the periodic curl of
$\lambda^2A(t,(x-c)/\lambda)$, so $v$ has zero spatial mean.

For the desired viscosity prescribe first
\[
 g=\partial_t v+(v\cdot\nabla)v-\nu\Delta v.
\]
This recomputes the residual at the torus viscosity; no invariance of
viscosity under spatial scaling is being assumed.  All mixed derivatives
of $g$ decrease rapidly, and its spatial mean is zero because $v$ has
mean zero, $(v\cdot\nabla)v=\operatorname{div}(v\otimes v)$, and
the Laplacian has zero integral.  Lemma~\ref{lem:projection} replaces $g$
effectively by its solenoidal, mean-zero projection, preserving all
these decay bounds and changing only the pressure, not the velocity.

In the indicated chart the particle trajectory is exactly
$c+\lambda X^W(t,0)$ as long as this expression remains in the chart,
by differentiation and uniqueness.  It remains there throughout every
nonhalting run and through the first halting signal.  During loading
its first coordinate is at least $3/8$; subsequently it is exactly
$3/8$ at every nonterminal step and throughout that step.  A terminal
step ends at $3/8-b_n/8$, which belongs to $(1/4,3/8)$ since $0<b_n<1$.
The remaining coordinates stay in the same interior chart by
\eqref{eq:rapid-total-memory} and the displacement bounds.  There is
therefore no wrapping ambiguity before the event, and the claimed
equivalence follows.
\end{proof}

